\documentclass[11pt,a4paper]{article}
% Common preamble for the hint sheets (\input from each file)
\usepackage[utf8]{inputenc}
\usepackage[T1]{fontenc}
\usepackage{lmodern}
\usepackage{amsmath,amssymb,amsthm,mathtools}
\usepackage[a4paper,margin=2.3cm]{geometry}
\usepackage{enumitem}
\usepackage{xcolor}
\usepackage[hidelinks]{hyperref}
\renewcommand{\leq}{\leqslant}
\renewcommand{\le}{\leqslant}
\renewcommand{\geq}{\geqslant}
\renewcommand{\ge}{\geqslant}
\setlength{\parindent}{0pt}
\setlength{\parskip}{0.35em}
% \hintstitle{Sheet title}{Class work / Homework}
\newcommand{\hintstitle}[2]{\begin{center}{\LARGE\bfseries Hints: #1\par}\vspace{0.5em}{\large #2\par}\end{center}
\noindent\textit{The hints to each problem are cumulative. Read one, try again, and open the next only if you are still stuck.}\par\medskip\hrule\medskip}
\newcommand{\problem}[1]{\par\bigskip{\large\bfseries Problem #1}\par\nopagebreak\smallskip}
\newcommand{\subpart}[1]{\par\smallskip\textbf{(#1)}\par\nopagebreak}
\newcommand{\hint}[1]{\par\smallskip\noindent\textcolor{black!70}{\textbf{Hint #1.}}\enspace}
\newcommand{\note}[1]{\par\smallskip\noindent\textit{Note.} #1\par}
\providecommand{\lt}{<}
\providecommand{\gt}{>}
\providecommand{\tightlist}{\setlength{\itemsep}{0pt}\setlength{\parskip}{0pt}}
\emergencystretch=3em

\begin{document}
\hintstitle{Number theory -- 5}{Class work}

\problem{1}
\hint{1} Theorem 1: let $g$ be the least positive element of $\{ax+by\}$. Division with remainder shows that $g$ divides $a$ and $b$, so $g=\gcd(a,b)$. If $(x_0,y_0)$ and $(x,y)$ are two solutions, then $a(x-x_0)=-b(y-y_0)$. Divide by $\gcd(a,b)$ and use coprimality to get the progression.
\hint{2} For $17x+13y=2026$: one particular solution is $(x,y)=(6,148)$, so all solutions are $x=6+13t$, $y=148-17t$, $t\in\mathbb Z$.
\hint{3} McNugget: for every $n$ there is a unique $y\in\{0,\dots,a-1\}$ with $by\equiv n\pmod a$. Show that $n$ is representable with $x,y\ge0$ if and only if $by\le n$ for this $y$. The worst residue is $y=a-1$, which gives the largest non-representable number $b(a-1)-a$.
\hint{4} Counting: for $0\le n\le ab-a-b$, show that exactly one of $n$ and $ab-a-b-n$ is representable. (If both were, adding the representations would represent $ab-a-b$.) So half of the $ab-a-b+1$ numbers are non-representable, which gives $(a-1)(b-1)/2$.

\problem{2}
\hint{1} Follow the proof sketched in Theorem 2. Dirichlet gives infinitely many $(x,y)$ with $|x^2-dy^2|<2\sqrt d+1$. So some value $k$ occurs infinitely often, and among those solutions two have $x\equiv x'$, $y\equiv y'\pmod k$. Check that $(x+y\sqrt d)/(x'+y'\sqrt d)$ has integer coordinates and norm $1$, and that it is not $\pm1$.
\hint{2} All solutions: let $\varepsilon=x_1+y_1\sqrt d$ be the least solution $>1$. If $u>1$ is a solution with $\varepsilon^n\le u<\varepsilon^{n+1}$, then $u\varepsilon^{-n}$ is a solution in $[1,\varepsilon)$, so it equals $1$.
\hint{3} $\sqrt{13}=[3;\overline{1,1,1,1,6}]$ has odd period $5$. The convergents are $3,4,\frac72,\frac{11}3,\frac{18}5$, and $18^2-13\cdot5^2=-1$. Squaring gives the fundamental solution of the $+1$ equation: $(18+5\sqrt{13})^2=649+180\sqrt{13}$.
\hint{4} For $d=34$ the fundamental solution of $x^2-34y^2=1$ is $35+6\sqrt{34}$. If the $-1$ equation were solvable, its least solution $\eta$ would satisfy $\eta^2=35+6\sqrt{34}$, i.e.\ $x^2+34y^2=35$ and $xy=3$, which is impossible. (Note that congruences cannot prove this: the $-1$ equation is solvable modulo every $n$.)

\problem{3}
\hint{1} Reduce modulo $a$ and modulo $a+1$. You get $a\mid2002$. Also $a+1$ divides $2002$ if $n$ is even and $2000$ if $n$ is odd.
\hint{2} The only divisors $a>1$ of $2002=2\cdot7\cdot11\cdot13$ with $a+1\mid2000$ or $a+1\mid2002$ are $a=7$ (with $n$ odd) and $a=13$ (with $n$ even). The case $a=1$ is impossible.
\hint{3} For $a=7$, look modulo $3$: this forces $n$ to be even, a contradiction. For $a=13$ with $n$ even, look modulo $8$: $13^{n+1}\equiv5$ while $2001\equiv1$, so $14^n\not\equiv0\pmod 8$ and $n\le2$. The unique solution is $a=13$, $n=2$: $13^3-14^2=2001$.

\problem{4}
\hint{1} Take a solution of $x^4+y^4=z^2$ with $z$ minimal. Then $(x^2,y^2,z)$ is a primitive Pythagorean triple, so (with $y$ even) $x^2=m^2-n^2$, $y^2=2mn$, $z=m^2+n^2$, where $m,n$ are coprime and $m$ is odd.
\hint{2} Then $x^2+n^2=m^2$ is again a primitive triple, so $n=2rs$ and $m=r^2+s^2$. Now $y^2=4rs(r^2+s^2)$ with $r,s,r^2+s^2$ pairwise coprime, so each is a square: $r=u^2$, $s=v^2$, $r^2+s^2=w^2$. This gives $u^4+v^4=w^2$ with $w<z$. A solution of $x^4+y^4=z^4$ would give one with $z^2$ in place of $z$.
\hint{3} Pythagorean triples: $(x,y,z)=k(m^2-n^2,\,2mn,\,m^2+n^2)$ with $m>n>0$ coprime and of opposite parity (and $x,y$ possibly swapped). Derive this from $y^2=(z-x)(z+x)$ in the primitive case.
\hint{4} For $x^2+y^2=2z^2$, $x$ and $y$ have the same parity. Put $u=\frac{x+y}2$, $v=\frac{x-y}2$; then $u^2+v^2=z^2$. So the solutions are $x=u+v$, $y=u-v$ with $(u,v,z)$ a Pythagorean triple, allowing $v\le0$. Example: $7^2+17^2=2\cdot13^2$.

\problem{5}
\hint{1} First pin down $x$. An odd prime $p\mid x$ does not divide $xyz+2$, so $p^3\mid2012=4\cdot503$, which is impossible. If $4\mid x$, compare the powers of $2$ on both sides. So $x\in\{1,2\}$.
\hint{2} In both cases $y^3+z^3\equiv0\pmod{503}$. Since $3\nmid502$, cubing is a bijection mod $503$, so $503\mid y+z$. Write $y+z=503k$ and divide by $y+z$, using $y^2-yz+z^2=(y+z)^2-3yz$.
\hint{3} For $x=2$ the equation becomes $y^3+z^3=503(yz+1)$, i.e.\ $k(y^2-yz+z^2)=yz+1$. Show that $k\ge2$ is impossible. So $k=1$ and $(y-z)^2=1$, giving $(x,y,z)=(2,251,252)$.
\hint{4} For $x=1$ you get $yz=\frac{503^2k^3-8}{3k+4}$. The bound $yz\le(y+z)^2/4$ forces $k\le4$. Check $k=1,2,3,4$: only $k=1$ gives an integer, and then $(y-z)^2=503^2-4yz$ is not a square. Answer: only $(2,251,252)$.
\end{document}
